\input{estilo}
% Las tres semánticas de entrega como tabla: la misma entrada (A le manda un
% registro a B), lo que hace el emisor y cuántas veces lo procesa B. Exactly-
% once es imposible en la práctica; se consigue con retry y un ID para
% descartar los repetidos, que es lo que hace MillWheel.
\begin{document}
\begin{figura}\begin{tikzpicture}[
    mini/.style={caja, minimum width=6mm, minimum height=6mm, inner sep=0pt, font=\scriptsize},
    enc/.style={etiq, text=suave, anchor=west}]
  \node[enc] at (0,0.9) {semántica};
  \node[enc] at (3.4,0.9) {el emisor};
  \node[enc] at (11.6,0.9) {B lo procesa};
  \draw[line width=0.5pt, draw=suave] (0,0.6) -- (15.2,0.6);
  \foreach \r/\s/\e/\res in {0/at-most-once/sin retry/{0 o 1 vez},
                             1/at-least-once/con retry/{1 o más veces},
                             2/exactly-once/{con retry + ID}/{exactamente 1 vez}} {
    \begin{scope}[yshift=-\r*1.6cm]
      \node[etiq, font=\small, anchor=west] at (0,0) {\s};
      \node[etiq, anchor=west] at (3.4,0) {\e};
      \node[mini] (a) at (6.6,0) {A};
      \node[mini] (b) at (9.6,0) {B};
    \end{scope}
  }
  % at-most-once: se manda una vez y se puede perder
  \draw[msg] (6.9,0) -- (8.1,0); \pic[scale=0.7] at (8.25,0) {rechazo};
  % at-least-once: el ack se pierde, A reintenta y B lo procesa dos veces
  \draw[msg] (6.9,-1.35) -- (9.3,-1.35);
  \draw[respuesta] (9.3,-1.6) -- (8.1,-1.6); \pic[scale=0.7] at (7.95,-1.6) {rechazo};
  \draw[msg] (6.9,-1.85) -- node[etiq, font=\scriptsize, below] {retry} (9.3,-1.85);
  % exactly-once: lo mismo, pero B reconoce el ID y descarta el repetido
  \draw[msg] (6.9,-2.95) -- (9.3,-2.95);
  \draw[respuesta] (9.3,-3.2) -- (8.1,-3.2); \pic[scale=0.7] at (7.95,-3.2) {rechazo};
  \draw[msg] (6.9,-3.45) -- node[etiq, font=\scriptsize, below] {retry, mismo ID} (9.3,-3.45);
  \node[etiq, text=suave, anchor=west] at (9.95,-3.45) {descarta};
  \foreach \r/\res in {0/{0 o 1 vez}, 1/{1 o más veces}, 2/{exactamente 1 vez}}
    \node[etiq, font=\small, anchor=west] at (11.6,-\r*1.6) {\res};
  \node[etiq, text=acento, anchor=west, align=left] at (0,-3.65) {como entrega,\\imposible en la práctica};
\end{tikzpicture}\end{figura}
\end{document}
